Geben Sie ein Beispiel für Nutzenfunktionen

Den Haushalten steht somit eine Menge M an zulässigen Kombinationen verschiedener Tee \( T \) und Kaffeemengen \( K \) zur Auswahl, es gilt: 
\( M \subset T \times K \).
Für diese Menge \( M \) wird in der ökonomischen Theorie jeweils eine Ordnungsrelation \( P \) mit \( P \subset T \times K \) 
(„Präferenzrelation“) definiert, diese erfüllt die
Anforderungen:
• Monotonie (11 Tassen Tee und 12 Tassen Kaffee werden immer 10 Tassen Tee und 11 Tassen Kaffee vorgezogen)
• Transitivität (Wenn 20 Tassen Tee und 10 Tassen Kaffee 30 Tassen Tee und 5 Tassen Kaffee vorgezogen werden, dann werden auch 20 Tassen
Tee und 30 Tassen Kaffee 25 Tassen Tee und 5 Tassen Kaffee vorgezogen)
Um die Anwendbarkeit der Relation zu vereinfachen wird eine sog. Nutzenfunktion \( U(T, K) \)
Es gilt in der Mikroökonomie das Prinzip des abnehmenden

Grenznutzen \( MU_K \)

Ein Haushalt wird in seiner Abwägung zwischen zwei
Konsumalternativen ein optimales Konsumbündel (aus der Menge der
zulässigen Kombinationen (z.B. Tee und Kaffee) und von Randlösungen
abstrahiert) die Alternative wählen, für die gilt:

Grenznutzen, d.h. der jeweilige Nutzenzuwachs einer neuen Tasse
nimmt ab, je mehr Tassen Kaffee.

\( MU_K = M U_T \)



Die individuelle Nachfragefunktion
Die Nachfragefunktion ist eine verkürzte Darstellung der optimalen
Ausgabeverteilung.
Preis
Für die Nachfrage nach einem Gut sind zwei Effekte besonders relevant.
Substitutionseffekt
Steigt der Preis für ein betrachtetes Gut steigen die Konsumente auf
alternative Produkte um. Steigt beispielsweise der Preis für Butter,
werden mehr Konsumenten Margarine kaufen.
Einkommenseffekt
Steigt der Preis eines Gutes, so wird es für die Konsumenten teurer,
sodass insgesamt weniger konsumiert werden.
• Normales Gut: Bei einer rückläufigen Kaufkraft geht die Nachfrage
nach dem Gut zurück
• Inferiores Gut: Bei einer rückläufigen Kaufkraft steigt die Nachfrage
nach dem Gut


SOZIALE WOLFAHRT

MORAL HAZARD & MARKET OF LEMONS
SIGNALLING
